\section{Introduction}
\subsection{Isoperimetric Inequalities Over the Hypercube}
Isoperimetric inequalities are fundamental results in mathematics with a myriad of applications. The main topic of this paper is discrete isoperimetric
inequalities, and more specifically isoperimetric inequalities over the hypercube $\{0,1\}^n$. Classical results along
these lines are due to Harper, Bernstein, Lindsey and Hart~\cite{Harper2,Bernstein,Lindsey,Hart}, Harper~\cite{Harper2}, Margulis~\cite{Margulis}, Talagrand~\cite{Tal1,Tal2,Tal4} and Kahn, Kalai and Linial~\cite{KKL}.

All of these theorems discuss various measures of boundaries and relations between them for Boolean functions.
For a function $f\colon\power{n}\to\power{}$ and a point $x\in\power{n}$, we define the \emph{sensitivity} of $f$ at $x$, denoted by $s_f(x)$,
to be the number of coordinates $i\in[n]$ such that $f(x\oplus e_i)\neq f(x)$. Thus, the \emph{vertex boundary} $V_f$ of a function $f$ is defined to be the set
of all $x$'s such that $s_f(x) > 0$, and the \emph{edge boundary} $E_f$ of $f$ is defined to be the set of edges $(x,x\oplus e_i)$ of the hypercube whose
endpoints are assigned different values by $f$. The most basic isoperimetric result related to Boolean functions,
known as Poincar\'{e}'s inequality, asserts that $\frac{\card{E_f}}{2^n}\geq {\sf var}(f)$,
i.e.\ the edge boundary of a function $f$ cannot be too small.\footnote{We remark that the term ``Poincar\'{e}'s inequality'' comes from the theory of functional inequalities
and Markov chains and refers to a much more general result, see for example~\cite[Section 2]{kitsos2014inequalities}. In the case of graphs, it asserts that for a graph
$G$, if the first non-trivial eigenvalue of its Laplacian $L$ is at least $\lambda$, then $\norm{L f}_2\geq \lambda \norm{f-\E[f]}_2$. In our
example, the Laplacian $L$ is $L = I-\frac{1}{n} A$ where $A$ is the adjacency matrix of the Boolean hypercube, and $\lambda$ is $2/n$.}
The edge isoperimetric inequality for the hypercube, due to Harper, Bernstein, Lindsey and Hart~\cite{Harper2,Bernstein,Lindsey,Hart}, is a strengthening
of that statement, asserting that for a fixed value of $\E[f]$, the functions that minimize the edge boundary $\frac{\card{E_f}}{2^n}$
are indicator functions
of initial segments in lexicographical orderings of the hypercube. In particular, this result implies that if $f$ is a non-constant Boolean function, then $\frac{\card{E_f}}{2^n}\geq \E[f]\log\left(\frac{1}{\E[f]}\right)$
(here and throughout, $\log$ means the logarithm function with base $2$). The vertex isoperimetric inequality, due to Harper~\cite{Harper2}, asserts that among all Boolean function
of a given measure, the functions that minimize the size of the vertex boundary $V_f$ are the indicator functions of
initial segments in simplicial orderings of the hypercube.
